% i18n_qa_conflicting_paper.tex
% Paper with conflicting babel options: last option is "english" → main=English
% BUT the body text is mostly French (deliberate mismatch)
% Expected detect_language result: "en" (last babel option wins)
% Deliberate errors: both French and English typography errors
\documentclass[11pt,a4paper]{article}
\usepackage[utf8]{inputenc}
\usepackage[T1]{fontenc}
\usepackage[french,english]{babel}
\usepackage{amsmath,amssymb,amsthm}
\usepackage{graphicx}
\usepackage{booktabs}
\usepackage{hyperref}
\usepackage{tikz}
\usepackage{pgfplots}
\pgfplotsset{compat=1.18}
\usepackage{float}

\newtheorem{theorem}{Theorem}[section]
\newtheorem{lemma}[theorem]{Lemma}
\newtheorem{proposition}[theorem]{Proposition}
\theoremstyle{definition}
\newtheorem{definition}{Definition}[section]

\title{Mod\`eles de transport optimal pour la segmentation\\
  d'images m\'edicales par r\'eseaux de neurones}
\author{Pierre-Antoine Moreau\textsuperscript{1}\quad
  Sophie Laurent\textsuperscript{2}\\[6pt]
  \textsuperscript{1}INRIA Paris, \'Equipe SIERRA\\
  \textsuperscript{2}\'Ecole Normale Sup\'erieure, D\'epartement d'Informatique}
\date{F\'evrier 2026}

\begin{document}
\maketitle

\begin{abstract}
% Abstract in English (matches babel main language)
We present a novel approach to medical image segmentation that leverages
optimal transport theory to define a geometrically meaningful loss function
for neural network training. Our method, termed Wasserstein Segmentation
Network (WSegNet), replaces the standard cross-entropy loss with a
regularized optimal transport distance between predicted and ground-truth
segmentation masks. We prove that the Sinkhorn divergence provides a
smooth approximation to the Wasserstein distance that is amenable to
backpropagation, and we establish convergence guarantees for the
resulting training procedure. Experiments on the ACDC cardiac MRI
dataset and the BraTS brain tumor segmentation benchmark demonstrate
improvements of 2.1\% and 1.8\% in Dice score respectively over
U-Net baselines trained with cross-entropy loss.

\medskip
\noindent\textbf{Keywords:} optimal transport, Wasserstein distance,
image segmentation, neural networks, medical imaging.
\end{abstract}


\section{Introduction}

% Body text mostly in French despite babel declaring English as main
% ERROR: missing NBSP before French punctuation (colon)
La segmentation d'images m\'edicales repr\'esente un enjeu majeur en
imagerie num\'erique. Les m\'ethodes fond\'ees sur les r\'eseaux de
neurones convolutifs, en particulier l'architecture U-Net propos\'ee par
Ronneberger et al.~\cite{ronneberger2015unet}, ont consid\'erablement
am\'elior\'e les performances de segmentation automatique. Cependant,
la fonction de perte utilis\'ee pour l'entra\^inement joue un r\^ole
d\'eterminant dans la qualit\'e des r\'esultats obtenus.

% ERROR: English-style double quotes in French text
La plupart des approches actuelles utilisent l'entropie crois\'ee ou le
coefficient de Dice comme fonction de perte. Ces fonctions traitent chaque
pixel "ind\'ependamment" et ignorent la structure g\'eom\'etrique de
l'espace des segmentations.

% ERROR: missing NBSP before exclamation mark
Dans ce travail, nous proposons d'utiliser la distance de Wasserstein
comme fonction de perte alternative! Cette distance, issue de la
th\'eorie du transport optimal, poss\`ede des propri\'et\'es
g\'eom\'etriques remarquables qui la rendent particuli\`erement
adapt\'ee \`a la comparaison de distributions spatiales.

\begin{definition}[Distance de Wasserstein]
Soient $\mu, \nu$ deux mesures de probabilit\'e sur $\mathbb{R}^d$.
La distance de Wasserstein d'ordre $p \geq 1$ est d\'efinie par:
\begin{equation}\label{eq:wasserstein}
  W_p(\mu, \nu) = \left( \inf_{\gamma \in \Pi(\mu, \nu)}
    \int_{\mathbb{R}^d \times \mathbb{R}^d} \|x - y\|^p \,
    \mathrm{d}\gamma(x, y) \right)^{1/p},
\end{equation}
o\`u $\Pi(\mu, \nu)$ d\'esigne l'ensemble des couplages de $\mu$ et $\nu$.
\end{definition}


\section{Transport optimal r\'egularis\'e}

% ERROR: missing NBSP before semicolon
Le calcul direct de la distance de Wasserstein est co\^uteux;
sa complexit\'e est $O(n^3 \log n)$ pour des mesures discr\`etes
\`a $n$ supports. La r\'egularisation entropique, introduite par
Cuturi~\cite{cuturi2013sinkhorn}, permet de r\'eduire cette
complexit\'e \`a $O(n^2 / \epsilon)$ o\`u $\epsilon$ est le
param\`etre de r\'egularisation.

\begin{definition}[Distance de Sinkhorn]
La distance de Sinkhorn r\'egularis\'ee est d\'efinie par:
\begin{equation}\label{eq:sinkhorn}
  S_\epsilon(\mu, \nu) = \text{OT}_\epsilon(\mu, \nu)
    - \frac{1}{2} \text{OT}_\epsilon(\mu, \mu)
    - \frac{1}{2} \text{OT}_\epsilon(\nu, \nu),
\end{equation}
o\`u
\begin{equation}\label{eq:ot-reg}
  \text{OT}_\epsilon(\mu, \nu) = \min_{\gamma \in \Pi(\mu, \nu)}
    \left\{ \int \|x-y\|^2 \, \mathrm{d}\gamma + \epsilon \,
    \text{KL}(\gamma \| \mu \otimes \nu) \right\}.
\end{equation}
\end{definition}

\begin{lemma}[Propri\'et\'es de diff\'erentiabilit\'e]
La divergence de Sinkhorn $S_\epsilon$ est:
\begin{enumerate}
  \item continument diff\'erentiable par rapport aux mesures d'entr\'ee;
  \item d\'efinie positive ($S_\epsilon(\mu, \nu) \geq 0$ avec \'egalit\'e
    ssi $\mu = \nu$);
  \item $|\sqrt{S_\epsilon(\mu, \nu)} - W_2(\mu, \nu)| = O(\sqrt{\epsilon})$
    quand $\epsilon \to 0$.
\end{enumerate}
\end{lemma}


\section{Architecture WSegNet}

L'architecture WSegNet repose sur un encodeur-d\'ecodeur de type U-Net
avec des connexions r\'esiduelles. La sortie du r\'eseau est une carte
de probabilit\'es $\hat{p}(x) \in [0, 1]^C$ pour chaque pixel $x$,
o\`u $C$ est le nombre de classes.

% ERROR: straight English quotes used in French text
La fonction de perte totale combine la divergence de Sinkhorn avec un
terme de r\'egularisation spatiale, formant ce que nous appelons la
"perte de transport":
\begin{equation}\label{eq:total-loss}
  \mathcal{L}(\theta) = \sum_{c=1}^C S_\epsilon(\hat{p}_c, p_c^*)
    + \lambda \sum_{c=1}^C \int \|\nabla \hat{p}_c(x)\|^2 \, \mathrm{d}x,
\end{equation}
o\`u $p_c^*$ est la segmentation de r\'ef\'erence pour la classe $c$
et $\lambda > 0$ contr\^ole la r\'egularisation.

\begin{theorem}[Convergence de l'entra\^inement]\label{thm:convergence}
Sous les hypoth\`eses standard de lissit\'e du r\'eseau et de bornitude
des gradients, l'algorithme SGD appliqu\'e \`a la perte~\eqref{eq:total-loss}
converge \`a un point stationnaire au taux:
\begin{equation}
  \min_{t \leq T} \mathbb{E}[\|\nabla_\theta \mathcal{L}(\theta_t)\|^2]
    \leq \frac{C_0}{\sqrt{T}},
\end{equation}
o\`u $C_0$ d\'epend de la constante de Lipschitz du gradient et de la
variance du bruit stochastique.
\end{theorem}


\section{Exp\'eriences}

\subsection{Protocole exp\'erimental}

% ERROR: double space after period (English style error)
Nous \'evaluons WSegNet sur deux benchmarks publics de segmentation
d'images m\'edicales.  Le premier est le jeu de donn\'ees ACDC
(Automated Cardiac Diagnosis Challenge) contenant 150 examens IRM
cardiaques avec segmentation du ventricule gauche, du ventricule droit
et du myocarde.

Le second benchmark est BraTS (Brain Tumor Segmentation), comprenant
335 volumes IRM c\'er\'ebraux avec annotation de trois types de tumeurs.

\begin{table}[H]
\centering
\caption{Description des jeux de donn\'ees}\label{tab:datasets}
\begin{tabular}{@{}lrrrr@{}}
\toprule
Benchmark & Sujets & Classes & R\'esolution & Modalit\'e \\
\midrule
ACDC & 150 & 4 & $256 \times 256$ & IRM cine \\
BraTS & 335 & 4 & $240 \times 240 \times 155$ & IRM T1/T2/FLAIR \\
\bottomrule
\end{tabular}
\end{table}

\subsection{R\'esultats quantitatifs}

% ERROR: missing NBSP before colon
Les r\'esultats principaux sont pr\'esent\'es dans le tableau~\ref{tab:results}.
WSegNet am\'eliore les performances sur les deux benchmarks:

\begin{table}[H]
\centering
\caption{Score de Dice (\%) sur les benchmarks ACDC et BraTS}\label{tab:results}
\begin{tabular}{@{}lcccc@{}}
\toprule
M\'ethode & \multicolumn{2}{c}{ACDC} & \multicolumn{2}{c}{BraTS} \\
\cmidrule(lr){2-3} \cmidrule(lr){4-5}
 & Dice & HD95 & Dice & HD95 \\
\midrule
U-Net + CE & 88.4 & 12.3 & 84.2 & 15.7 \\
U-Net + Dice & 89.1 & 11.8 & 84.9 & 14.9 \\
Attention U-Net & 89.7 & 10.5 & 85.6 & 13.8 \\
TransUNet & 90.2 & 9.8 & 86.1 & 13.2 \\
\textbf{WSegNet (ours)} & \textbf{90.5} & \textbf{8.9} & \textbf{86.0} & \textbf{12.4} \\
\bottomrule
\end{tabular}
\end{table}

On observe que WSegNet obtient les meilleurs scores de Dice sur ACDC
(90.5\%) et des performances comp\'etitives sur BraTS (86.0\%). La
m\'etrique HD95 (distance de Hausdorff au 95e percentile) est
am\'elior\'ee de mani\`ere significative, ce qui confirme que la perte
de transport favorise des contours de segmentation plus pr\'ecis.

\subsection{\'Etude d'ablation}

\begin{table}[H]
\centering
\caption{\'Etude d'ablation sur ACDC}\label{tab:ablation}
\begin{tabular}{@{}lccc@{}}
\toprule
Configuration & Dice (\%) & HD95 (mm) & Temps/\'epoque (s) \\
\midrule
CE seule & 88.4 & 12.3 & 45 \\
Sinkhorn seule & 89.8 & 9.5 & 78 \\
Sinkhorn + TV & 90.1 & 9.2 & 82 \\
Sinkhorn + $H^1$ reg. & \textbf{90.5} & \textbf{8.9} & 85 \\
\bottomrule
\end{tabular}
\end{table}


\section{Analyse th\'eorique}

\begin{proposition}[Avantage g\'eom\'etrique du transport optimal]
Pour deux masques binaires $A, B \subset \Omega$ tels que $B$ est une
translation de $A$ par un vecteur $v$, on a:
\begin{equation}
  W_2(\mathbf{1}_A, \mathbf{1}_B) = \|v\| \cdot \sqrt{|A|},
\end{equation}
tandis que l'entropie crois\'ee ne d\'epend que du recouvrement:
\begin{equation}
  \text{CE}(\mathbf{1}_A, \mathbf{1}_B) = -|A \cap B| \log 1 -
    |A \setminus B| \log 0 = +\infty.
\end{equation}
Ceci illustre que $W_2$ fournit un signal de gradient informatif m\^eme
lorsque les masques ne se recouvrent pas, contrairement \`a l'entropie
crois\'ee.
\end{proposition}

\begin{figure}[H]
\centering
\begin{tikzpicture}
\begin{axis}[
  width=0.75\textwidth,
  height=5.5cm,
  xlabel={\'Epoque},
  ylabel={Score de Dice (\%)},
  legend pos=south east,
  grid=major,
  xmin=0, xmax=200,
  ymin=70, ymax=92,
]
\addplot[blue, thick] coordinates {
  (0,72) (20,80) (40,84) (60,86.5) (80,88) (100,89)
  (120,89.5) (140,90) (160,90.3) (180,90.4) (200,90.5)
};
\addplot[red, thick, dashed] coordinates {
  (0,71) (20,78) (40,82) (60,84.5) (80,86) (100,87)
  (120,87.5) (140,88) (160,88.2) (180,88.3) (200,88.4)
};
\addplot[green!60!black, thick, dotted] coordinates {
  (0,72.5) (20,79.5) (40,83.5) (60,85.5) (80,87.5) (100,88.5)
  (120,89) (140,89.3) (160,89.5) (180,89.6) (200,89.7)
};
\legend{WSegNet, U-Net+CE, Attention U-Net}
\end{axis}
\end{tikzpicture}
\caption{Courbes d'entra\^inement sur ACDC (validation Dice).
  WSegNet converge plus rapidement et atteint un meilleur score
  final.}\label{fig:training-curves}
\end{figure}


\section{Conclusion}

% Mixed French/English conclusion
Nous avons pr\'esent\'e WSegNet, une approche de segmentation d'images
m\'edicales fond\'ee sur le transport optimal. Les r\'esultats
exp\'erimentaux d\'emontrent l'int\'er\^et de la divergence de Sinkhorn
comme fonction de perte pour l'entra\^inement de r\'eseaux de neurones
de segmentation. Les perspectives incluent l'extension aux probl\`emes
3D volum\'etriques et l'\'etude de variantes non r\'egularis\'ees
utilisant des formulations duales.

% ERROR: missing NBSP before question mark (French style error)
Une question ouverte importante reste celle de la g\'en\'eralisation
multi-classes: comment \'etendre efficacement la perte de transport
au cas o\`u le nombre de classes est sup\'erieur \`a deux?


\begin{thebibliography}{99}

\bibitem{ronneberger2015unet}
O.~Ronneberger, P.~Fischer, and T.~Brox.
\newblock {U-Net}: Convolutional networks for biomedical image segmentation.
\newblock In \textit{Medical Image Computing and Computer-Assisted Intervention
  (MICCAI)}, pages 234--241, 2015.

\bibitem{cuturi2013sinkhorn}
M.~Cuturi.
\newblock Sinkhorn distances: Lightspeed computation of optimal transport.
\newblock In \textit{Advances in Neural Information Processing Systems
  (NeurIPS)}, pages 2292--2300, 2013.

\bibitem{peyre2019computational}
G.~Peyr\'e and M.~Cuturi.
\newblock Computational optimal transport: With applications to data science.
\newblock \textit{Foundations and Trends in Machine Learning},
  11(5--6):355--607, 2019.

\bibitem{feydy2019interpolating}
J.~Feydy, T.~S\'ejourn\'e, F.-X.~Vialard, S.-I.~Amari, A.~Trouv\'e,
  and G.~Peyr\'e.
\newblock Interpolating between optimal transport and {MMD} using {Sinkhorn}
  divergences.
\newblock In \textit{International Conference on Artificial Intelligence and
  Statistics (AISTATS)}, pages 2681--2690, 2019.

\bibitem{villani2009optimal}
C.~Villani.
\newblock \textit{Optimal Transport: Old and New}.
\newblock Springer, 2009.

\bibitem{chen2018transunet}
J.~Chen, Y.~Lu, Q.~Yu, et~al.
\newblock {TransUNet}: Transformers make strong encoders for medical image
  segmentation.
\newblock \textit{arXiv preprint arXiv:2102.04306}, 2021.

\end{thebibliography}

\end{document}
